%% netlayer.tex
%%
%% Copyright (C) 2000-2018 Simone Piccardi.  Permission is granted to
%% copy, distribute and/or modify this document under the terms of the GNU Free
%% Documentation License, Version 1.1 or any later version published by the
%% Free Software Foundation; with the Invariant Sections being "Un preambolo",
%% with no Front-Cover Texts, and with no Back-Cover Texts.  A copy of the
%% license is included in the section entitled "GNU Free Documentation
%% License".
%% 

\chapter{Il livello di rete}
\label{cha:network_layer}

In questa appendice prenderemo in esame i vari protocolli disponibili a
livello di rete.\footnote{per la spiegazione della suddivisione in livelli dei
  protocolli di rete, si faccia riferimento a quanto illustrato in
  sez.~\ref{sec:net_protocols}.} Per ciascuno di essi forniremo una descrizione
generica delle principali caratteristiche, del formato di dati usato e quanto
possa essere necessario per capirne meglio il funzionamento dal punto di vista
della programmazione.

Data la loro prevalenza il capitolo sarà sostanzialmente incentrato sui due
protocolli principali esistenti su questo livello: il protocollo IP, sigla che
sta per \textit{Internet Protocol}, (ma che più propriamente si dovrebbe
chiamare IPv4) ed la nuova versione di questo stesso protocollo, denominata
IPv6. Tratteremo comunque anche il protocollo ICMP e la sua versione
modificata per IPv6 (cioè ICMPv6).


\section{Il protocollo IP}
\label{sec:ip_protocol}

L'attuale \textit{Internet Protocol} (IPv4) viene standardizzato nel 1981
dall'\href{http://www.ietf.org/rfc/rfc791.txt}{RFC~791}; esso nasce per
disaccoppiare le applicazioni della struttura hardware delle reti di
trasmissione, e creare una interfaccia di trasmissione dei dati indipendente
dal sottostante substrato di rete, che può essere realizzato con le tecnologie
più disparate (Ethernet, Token Ring, FDDI, ecc.).


\subsection{Introduzione}
\label{sec:IP_intro}

Il compito principale di IP è quello di trasmettere i pacchetti da un computer
all'altro della rete; le caratteristiche essenziali con cui questo viene
realizzato in IPv4 sono due:
\begin{itemize}
\item \textit{Universal addressing} la comunicazione avviene fra due host
  identificati univocamente con un indirizzo a 32 bit che può appartenere ad
  una sola interfaccia di rete.
\item \textit{Best effort} viene assicurato il massimo impegno nella
  trasmissione, ma non c'è nessuna garanzia per i livelli superiori né
  sulla percentuale di successo né sul tempo di consegna dei pacchetti di
  dati, né sull'ordine in cui vengono consegnati.
\end{itemize}

Per effettuare la comunicazione e l'instradamento dei pacchetti fra le varie
reti di cui è composta Internet IPv4 organizza gli indirizzi in una
gerarchia a due livelli, in cui una parte dei 32 bit dell'indirizzo indica il
numero di rete, e un'altra l'host al suo interno.  Il numero di rete serve
ai router per stabilire a quale rete il pacchetto deve essere inviato, il
numero di host indica la macchina di destinazione finale all'interno di detta
rete.

Per garantire l'unicità dell'indirizzo Internet esiste un'autorità centrale
(la IANA, \textit{Internet Assigned Number Authority}) che assegna i numeri di
rete alle organizzazioni che ne fanno richiesta; è poi compito di quest'ultime
assegnare i numeri dei singoli host all'interno della propria rete.

Per venire incontro alle richieste dei vari enti e organizzazioni che volevano
utilizzare questo protocollo di comunicazione, originariamente gli indirizzi
di rete erano stati suddivisi all'interno delle cosiddette \textit{classi},
(rappresentate in tab.~\ref{tab:IP_ipv4class}), in modo da consentire
dispiegamenti di reti di varie dimensioni a seconda delle diverse esigenze.

\begin{table}[htb]
  \centering
  \footnotesize
  \begin{usepicture} 
  \begin{tabular} {c@{\hspace{1mm}\vrule}
      p{3mm}@{\vrule}p{3mm}@{\vrule}p{3mm}@{\vrule}p{3mm}@{\vrule}
      p{3mm}@{\vrule}p{3mm}@{\vrule}p{3mm}@{\vrule}p{3mm}@{\vrule}
      p{3mm}@{\vrule}p{3mm}@{\vrule}p{3mm}@{\vrule}p{3mm}@{\vrule}
      p{3mm}@{\vrule}p{3mm}@{\vrule}p{3mm}@{\vrule}p{3mm}@{\vrule}
      p{3mm}@{\vrule}p{3mm}@{\vrule}p{3mm}@{\vrule}p{3mm}@{\vrule}
      p{3mm}@{\vrule}p{3mm}@{\vrule}p{3mm}@{\vrule}p{3mm}@{\vrule}
      p{3mm}@{\vrule}p{3mm}@{\vrule}p{3mm}@{\vrule}p{3mm}@{\vrule}
      p{3mm}@{\vrule}p{3mm}@{\vrule}p{3mm}@{\vrule}p{3mm}@{\vrule}}
    \omit&\omit& \multicolumn{7}{c}{7 bit}&\multicolumn{24}{c}{24 bit} \\
    \cline{2-33}
    \omit\hfill\vrule &&&&&&&& &&&&&&&& &&&&&&&& &&&&&&&& \\
    classe A &\centering 0&
    \multicolumn{7}{@{}c@{\vrule}}{\parbox[c]{21mm}{\centering net Id}} &
    \multicolumn{24}{@{}c@{\vrule}}{\parbox[c]{72mm}{\centering host Id}} \\
    \omit\hfill\vrule &&&&&&&& &&&&&&&& &&&&&&&& &&&&&&&& \\
    \cline{2-33}
    \multicolumn{33}{c}{ } \\
    \omit&\omit&\omit& 
    \multicolumn{14}{c}{14 bit}&\multicolumn{16}{c}{16 bit} \\
    \cline{2-33}
    \omit\hfill\vrule &&&&&&&& &&&&&&&& &&&&&&&& &&&&&&&& \\
    classe B&\centering 1&\centering 0& 
    \multicolumn{14}{@{}c@{\vrule}}{\parbox{42mm}{\centering net Id}} &
    \multicolumn{16}{@{}c@{\vrule}}{\parbox{48mm}{\centering host Id}} \\
    \omit\hfill\vrule &&&&&&&& &&&&&&&& &&&&&&&& &&&&&&&& \\
    \cline{2-33}
   
    \multicolumn{33}{c}{ } \\
    \omit&\omit&\omit& 
    \multicolumn{21}{c}{21 bit}&\multicolumn{8}{c}{8 bit} \\
    \cline{2-33}
    \omit\hfill\vrule &&&&&&&& &&&&&&&& &&&&&&&& &&&&&&&& \\
    classe C&\centering 1&\centering 1&\centering 0&
    \multicolumn{21}{@{}c@{\vrule}}{\parbox{63mm}{\centering net Id}} &
    \multicolumn{8}{@{}c@{\vrule}}{\parbox{24mm}{\centering host Id}} \\
    \omit\hfill\vrule &&&&&&&& &&&&&&&& &&&&&&&& &&&&&&&& \\
    \cline{2-33}


    \multicolumn{33}{c}{ } \\
    \omit&\omit&\omit&\omit& 
    \multicolumn{28}{c}{28 bit} \\
    \cline{2-33}
    \omit\hfill\vrule &&&&&&&& &&&&&&&& &&&&&&&& &&&&&&&& \\
    classe D&\centering 1&\centering 1&\centering 1&\centering 0&
    \multicolumn{28}{@{}c@{\vrule}}{\parbox{63mm}{\centering 
        multicast group Id}} \\
    \omit\hfill\vrule &&&&&&&& &&&&&&&& &&&&&&&& &&&&&&&& \\
    \cline{2-33}

    \multicolumn{33}{c}{ } \\
    \omit&\omit&\omit&\omit&\omit&
    \multicolumn{27}{c}{27 bit} \\
    \cline{2-33}
    \omit\hfill\vrule &&&&&&&& &&&&&&&& &&&&&&&& &&&&&&&& \\
    classe E&\centering 1&\centering 1&\centering 1&\centering 1&\centering 0&
    \multicolumn{27}{@{}c@{\vrule}}{\parbox{59mm}{\centering 
        reserved for future use}} \\
    \omit\hfill\vrule &&&&&&&& &&&&&&&& &&&&&&&& &&&&&&&& \\
    \cline{2-33}

  \end{tabular}
  \end{usepicture} 
\caption{Le classi di indirizzi secondo IPv4.}
\label{tab:IP_ipv4class}
\end{table}

Le classi di indirizzi usate per il dispiegamento delle reti su quella che
comunemente viene chiamata \textit{Internet} sono le prime tre; la classe D è
destinata al \textit{multicast} mentre la classe E è riservata per usi
sperimentali e non viene impiegata.

Come si può notare però la suddivisione riportata in
tab.~\ref{tab:IP_ipv4class} è largamente inefficiente in quanto se ad un
utente necessita anche solo un indirizzo in più dei 256 disponibili con una
classe A occorre passare a una classe B, che ne prevede 65536,\footnote{in
  realtà i valori esatti sarebbero 254 e 65534, una rete con a disposizione
  $N$ bit dell'indirizzo IP, ha disponibili per le singole macchine soltanto
  $@^N-2$ numeri, dato che uno deve essere utilizzato come indirizzo di rete e
  uno per l'indirizzo di \textit{broadcast}.} con un conseguente spreco di
numeri.

Inoltre, in particolare per le reti di classe C, la presenza di tanti
indirizzi di rete diversi comporta una crescita enorme delle tabelle di
instradamento che ciascun router dovrebbe tenere in memoria per sapere dove
inviare il pacchetto, con conseguente crescita dei tempi di elaborazione da
parte di questi ultimi ed inefficienza nel trasporto.

\begin{table}[htb]
  \centering
  \footnotesize
  \begin{usepicture} 
  \begin{tabular} {c@{\hspace{1mm}\vrule}
      p{3mm}@{\vrule}p{3mm}@{\vrule}p{3mm}@{\vrule}p{3mm}@{\vrule}
      p{3mm}@{\vrule}p{3mm}@{\vrule}p{3mm}@{\vrule}p{3mm}@{\vrule}
      p{3mm}@{\vrule}p{3mm}@{\vrule}p{3mm}@{\vrule}p{3mm}@{\vrule}
      p{3mm}@{\vrule}p{3mm}@{\vrule}p{3mm}@{\vrule}p{3mm}@{\vrule}
      p{3mm}@{\vrule}p{3mm}@{\vrule}p{3mm}@{\vrule}p{3mm}@{\vrule}
      p{3mm}@{\vrule}p{3mm}@{\vrule}p{3mm}@{\vrule}p{3mm}@{\vrule}
      p{3mm}@{\vrule}p{3mm}@{\vrule}p{3mm}@{\vrule}p{3mm}@{\vrule}
      p{3mm}@{\vrule}p{3mm}@{\vrule}p{3mm}@{\vrule}p{3mm}@{\vrule}}
    \omit&
    \multicolumn{12}{c}{$n$ bit}&\multicolumn{20}{c}{$32-n$ bit} \\
    \cline{2-33}
    \omit\hfill\vrule &&&&&&&& &&&&&&&& &&&&&&&& &&&&&&&& \\
    CIDR &
    \multicolumn{12}{@{}c@{\vrule}}{\parbox[c]{36mm}{\centering net Id}} &
    \multicolumn{20}{@{}c@{\vrule}}{\parbox[c]{60mm}{\centering host Id}} \\
    \omit\hfill\vrule &&&&&&&& &&&&&&&& &&&&&&&& &&&&&&&& \\
    \cline{2-33}
  \end{tabular}
  \end{usepicture} 
\caption{Uno esempio di indirizzamento CIDR.}
\label{tab:IP_ipv4cidr}
\end{table}

Per questo nel 1992 è stato introdotto un indirizzamento senza classi (il
CIDR, \textit{Classless Inter-Domain Routing}) in cui il limite fra i bit
destinati a indicare il numero di rete e quello destinati a indicare l'host
finale può essere piazzato in qualunque punto (vedi
tab.~\ref{tab:IP_ipv4cidr}), permettendo di accorpare più classi A su un'unica
rete o suddividere una classe B e diminuendo al contempo il numero di
indirizzi di rete da inserire nelle tabelle di instradamento dei router.


\subsection{L'intestazione di IP}
\label{sec:IP_header}

Come illustrato in fig.~\ref{fig:net_tcpip_data_flux} (si ricordi quanto detto
in sez.~\ref{sec:net_tcpip_overview} riguardo al funzionamento generale del
TCP/IP), per eseguire il suo compito il protocollo IP inserisce (come
praticamente ogni protocollo di rete) una opportuna intestazione in cima ai
dati che deve trasmettere, la cui schematizzazione è riportata in
fig.~\ref{fig:IP_ipv4_head}.

\begin{figure}[!htb]
  \centering
  \includegraphics[width=10cm]{img/ipv4_head}
  \caption{L'intestazione o \textit{header} di IPv4.}
  \label{fig:IP_ipv4_head}
\end{figure}

Ciascuno dei campi illustrati in fig.~\ref{fig:IP_ipv4_head} ha un suo preciso
scopo e significato, che si è riportato brevemente in
tab.~\ref{tab:IP_ipv4field}; si noti come l'intestazione riporti sempre due
indirizzi IP, quello \textsl{sorgente}, che indica l'IP da cui è partito il
pacchetto (cioè l'indirizzo assegnato alla macchina che lo spedisce) e quello
\textsl{destinazione} che indica l'indirizzo a cui deve essere inviato il
pacchetto (cioè l'indirizzo assegnato alla macchina che lo riceverà).

\begin{table}[!htb]
  \footnotesize
  \begin{center}
    \begin{tabular}{|l|c|p{10cm}|}
      \hline
      \textbf{Nome} & \textbf{Bit} & \textbf{Significato} \\
      \hline
      \hline
      \textit{version}       & 4& Numero di \textsl{versione}, nel caso 
                                  specifico vale sempre 4.\\
      \textit{head length}   & 4& Lunghezza dell'intestazione,
                                  in multipli di 32 bit.\\
      \textit{type of service}&8& Il ``\textsl{tipo di servizio}'', è suddiviso
                                  in: 3 bit di precedenza, che nelle attuali
                                  implementazioni del protocollo non vengono
                                  comunque utilizzati; un bit riservato che
                                  deve essere mantenuto a 0; 4 bit che
                                  identificano il tipo di servizio
                                  richiesto, uno solo dei quali può essere
                                  attivo.\\ 
      \textit{total length}  &16& La \textsl{lunghezza totale}, indica 
                                  la dimensione del carico di dati del
                                  pacchetto IP in byte.\\ 
      \textit{identification}&16& L'\textsl{identificazione}, assegnato alla
                                  creazione, è aumentato di uno all'origine
                                  della trasmissione di ciascun pacchetto, ma
                                  resta lo stesso per i pacchetti
                                  frammentati, consentendo così di
                                  identificare quelli che derivano dallo
                                  stesso pacchetto originario.\\
      \textit{flag}          & 3& I \textsl{flag} di controllo nell'ordine: il 
                                  primo è riservato e sempre nullo, il secondo
                                  indica se il pacchetto non può essere 
                                  frammentato, il terzo se ci sono ulteriori
                                  frammenti.\\  
      \textit{fragmentation offset}&13& L'\textsl{offset di frammento}, indica
                                  la posizione del frammento rispetto al
                                  pacchetto originale.\\
      \textit{time to live}  &16& Il \textsl{tempo di vita}, è decrementato di
                                  uno ogni volta che un router ritrasmette il
                                  pacchetto, se arriva a zero il pacchetto
                                  viene scartato.\\ 
      \textit{protocol}      & 8& Il \textsl{protocollo}, identifica il tipo di
                                  pacchetto che segue l'intestazione di IPv4.\\
      \textit{header checksum}&16&La \textsl{checksum di intestazione}, somma 
                                  di controllo per l'intestazione.\\ 
      \textit{source IP}     &32& L'\textsl{indirizzo di origine}.\\
      \textit{destination IP}&32& L'\textsl{indirizzo di destinazione}.\\
      \hline
    \end{tabular}
    \caption{Legenda per il significato dei campi dell'intestazione di IPv4}
    \label{tab:IP_ipv4field}
  \end{center}
\end{table}


Il campo TOS definisce il cosiddetto \textit{Type of Service}; questo permette
di definire il tipo di traffico contenuto nei pacchetti, e può essere
utilizzato dai router per dare diverse priorità in base al valore assunto da
questo campo. Abbiamo già visto come il valore di questo campo può essere
impostato sul singolo socket con l'opzione \const{IP\_TOS} (vedi
sez.~\ref{sec:sock_ipv4_options}), esso inoltre può essere manipolato sia dal
sistema del \textit{netfilter} di Linux con il comando \texttt{iptables} che
dal sistema del routing avanzato del comando \texttt{ip route} per consentire
un controllo più dettagliato dell'instradamento dei pacchetti e l'uso di
priorità e politiche di distribuzione degli stessi.

\begin{table}[!htb]
  \centering
  \footnotesize
  \begin{tabular}{|l|l|p{8cm}|}
    \hline
    \multicolumn{2}{|c|}{\textbf{Valore}}&\textbf{Significato}\\
    \hline
    \hline
    \constd{IPTOS\_LOWDELAY}    &\texttt{0x10}& Minimizza i ritardi
                                                per rendere più veloce
                                                possibile la ritrasmissione
                                                dei pacchetti (usato per
                                                traffico interattivo di
                                                controllo come SSH).\\
    \constd{IPTOS\_THROUGHPUT}  &\texttt{0x8} & Ottimizza la trasmissione
                                                per rendere il più elevato
                                                possibile il flusso netto di
                                                dati (usato su traffico dati,
                                                come quello di FTP).\\ 
    \constd{IPTOS\_RELIABILITY} &\texttt{0x4} & Ottimizza la trasmissione
                                                per ridurre al massimo le
                                                perdite di pacchetti (usato su
                                                traffico soggetto a rischio di
                                                perdita di pacchetti come TFTP
                                                o DHCP).\\
    \constd{IPTOS\_MINCOST}     &\texttt{0x2} & Indica i dati di riempimento,
                                                dove non interessa se si ha
                                                una bassa velocità di
                                                trasmissione, da utilizzare
                                                per i collegamenti con minor
                                                costo (usato per i protocolli
                                                di streaming).\\ 
    \textit{Normal-Service}&\texttt{0x0}      & Nessuna richiesta specifica.\\
    \hline

    \hline
  \end{tabular}
  \caption{Le costanti che definiscono alcuni valori standard per il campo TOS
    da usare come argomento \param{optval} per l'opzione \const{IP\_TOS}.} 
  \label{tab:IP_TOS_values}
\end{table}

I possibili valori del campo TOS, insieme al relativo significato ed alle
costanti numeriche ad esso associati, sono riportati in
tab.~\ref{tab:IP_TOS_values}. Per il valore nullo, usato di default per tutti
i pacchetti, e relativo al traffico normale, non esiste nessuna costante
associata. 

Il campo TTL, acromino di \textit{Time To Live}, viene utilizzato per
stabilire una sorta di tempo di vita massimo dei pacchetti sulla rete. In
realtà più che di un tempo, il campo serve a limitare il numero massimo di
salti (i cosiddetti \textit{hop}) che un pacchetto IP può compiere nel passare
da un router ad un altro nel suo attraversamento della rete verso la
destinazione.

Il protocollo IP prevede infatti che il valore di questo campo venga
decrementato di uno da ciascun router che ritrasmette il pacchetto verso la
sua destinazione, e che quando questo diventa nullo il router lo debba
scartare, inviando all'indirizzo sorgente un pacchetto ICMP di tipo
\textit{time-exceeded} con un codice \textit{ttl-zero-during-transit} se
questo avviene durante il transito sulla rete o
\textit{ttl-zero-during-reassembly} se questo avviene alla destinazione finale
(vedi sez.~\ref{sec:ICMP_protocol}).

In sostanza grazie all'uso di questo accorgimento un pacchetto non può
continuare a vagare indefinitamente sulla rete, e viene comunque scartato dopo
un certo tempo, o meglio, dopo che ha attraversato in certo numero di
router. Nel caso di Linux il valore iniziale utilizzato normalmente è 64 (vedi
sez.~\ref{sec:sock_ipv4_sysctl}).



\subsection{Le opzioni di IP}
\label{sec:IP_options}

Da fare ...


\section{Il protocollo IPv6}
\label{sec:ipv6_protocol}

Negli anni '90 con la crescita del numero di macchine connesse su Internet si
arrivò a temere l'esaurimento dello spazio degli indirizzi disponibili, specie
in vista di una prospettiva (per ora rivelatasi prematura) in cui ogni
apparecchio elettronico sarebbe stato inserito all'interno della rete. 

Per questo motivo si iniziò a progettare una nuova versione del protocollo 

L'attuale Internet Protocol (IPv4) viene standardizzato nel 1981
dall'\href{http://www.ietf.org/rfc/rfc719.txt}{RFC~719}; esso nasce per
disaccoppiare le applicazioni della struttura hardware delle reti di
trasmissione, e creare una interfaccia di trasmissione dei dati indipendente
dal sottostante substrato di rete, che può essere realizzato con le tecnologie
più disparate (Ethernet, Token Ring, FDDI, ecc.).


\subsection{I motivi della transizione}
\label{sec:IP_whyipv6}

Negli ultimi anni la crescita vertiginosa del numero di macchine connesse a
internet ha iniziato a far emergere i vari limiti di IPv4; in particolare si
è iniziata a delineare la possibilità di arrivare a una carenza di
indirizzi disponibili.

In realtà il problema non è propriamente legato al numero di indirizzi
disponibili; infatti con 32 bit si hanno $2^{32}$, cioè circa 4 miliardi,
numeri diversi possibili, che sono molti di più dei computer attualmente
esistenti.

Il punto è che la suddivisione di questi numeri nei due livelli rete/host e
l'utilizzo delle classi di indirizzamento mostrate in precedenza, ha
comportato che, nella sua evoluzione storica, il dispiegamento delle reti e
l'allocazione degli indirizzi siano stati inefficienti; neanche l'uso del CIDR
ha permesso di eliminare le inefficienze che si erano formate, dato che il
ridispiegamento degli indirizzi comporta cambiamenti complessi a tutti i
livelli e la riassegnazione di tutti gli indirizzi dei computer di ogni
sottorete.

Diventava perciò necessario progettare un nuovo protocollo che permettesse
di risolvere questi problemi, e garantisse flessibilità sufficiente per
poter continuare a funzionare a lungo termine; in particolare necessitava un
nuovo schema di indirizzamento che potesse rispondere alle seguenti
necessità:

\begin{itemize}
\item un maggior numero di numeri disponibili che consentisse di non restare
  più a corto di indirizzi
\item un'organizzazione gerarchica più flessibile dell'attuale 
\item uno schema di assegnazione degli indirizzi in grado di minimizzare le
  dimensioni delle tabelle di instradamento
\item uno spazio di indirizzi che consentisse un passaggio automatico dalle
  reti locali a internet
\end{itemize}


\subsection{Principali caratteristiche di IPv6}
\label{sec:IP_ipv6over}

Per rispondere alle esigenze descritte in sez.~\ref{sec:IP_whyipv6} IPv6 nasce
come evoluzione di IPv4, mantenendone inalterate le funzioni che si sono
dimostrate valide, eliminando quelle inutili e aggiungendone poche altre
ponendo al contempo una grande attenzione a mantenere il protocollo il più
snello e veloce possibile.

I cambiamenti apportati sono comunque notevoli e possono essere riassunti a
grandi linee nei seguenti punti:
\begin{itemize}
\item l'espansione delle capacità di indirizzamento e instradamento, per
  supportare una gerarchia con più livelli di indirizzamento, un numero di
  nodi indirizzabili molto maggiore e una auto-configurazione degli indirizzi
\item l'introduzione un nuovo tipo di indirizzamento, l'\textit{anycast} che
  si aggiungono agli usuali \textit{unicast} e \textit{multicast}
\item la semplificazione del formato dell'intestazione, eliminando o rendendo
  opzionali alcuni dei campi di IPv4, per eliminare la necessità di
  riprocessare la stessa da parte dei router e contenere l'aumento di
  dimensione dovuto ai nuovi indirizzi
\item un supporto per le opzioni migliorato, per garantire una trasmissione
  più efficiente del traffico normale, limiti meno stringenti sulle
  dimensioni delle opzioni, e la flessibilità necessaria per introdurne di
  nuove in futuro
\item il supporto per delle capacità di qualità di servizio (QoS) che
  permetta di identificare gruppi di dati per i quali si può provvedere un
  trattamento speciale (in vista dell'uso di internet per applicazioni
  multimediali e/o ``real-time'')
\end{itemize}


\subsection{L'intestazione di IPv6}
\label{sec:IP_ipv6head}

Per capire le caratteristiche di IPv6 partiamo dall'intestazione usata dal
protocollo per gestire la trasmissione dei pacchetti; in
fig.~\ref{fig:IP_ipv6head} è riportato il formato dell'intestazione di IPv6 da
confrontare con quella di IPv4 in fig.~\ref{fig:IP_ipv4_head}. La spiegazione
del significato dei vari campi delle due intestazioni è riportato
rispettivamente in tab.~\ref{tab:IP_ipv6field} e tab.~\ref{tab:IP_ipv4field})

% \begin{table}[htb]
%   \footnotesize
%   \begin{center}
%     \begin{tabular}{@{\vrule}p{16mm}@{\vrule}p{16mm}@{\vrule}p{16mm}
%         @{\vrule}p{16mm}@{\vrule}p{16mm}@{\vrule}p{16mm}
%         @{\vrule}p{16mm}@{\vrule}p{16mm}@{\vrule} }
%     \multicolumn{8}{@{}c@{}}{0\hfill 15 16\hfill 31}\\
%     \hline
%     \centering version&\centering priority& 
%     \multicolumn{6}{@{}p{96mm}@{\vrule}}{\centering flow label} \\
%     \hline
%     \multicolumn{4}{@{\vrule}p{64mm}@{\vrule}}{\centering payload length} & 
%     \multicolumn{2}{@{}p{32mm}@{\vrule}}{\centering next header} & 
%     \multicolumn{2}{@{}p{32mm}@{\vrule}}{\centering hop limit}\\
%     \hline
%     \multicolumn{8}{@{\vrule}c@{\vrule}}{} \\
%     \multicolumn{8}{@{\vrule}c@{\vrule}}{
%       source} \\
%     \multicolumn{8}{@{\vrule}c@{\vrule}}{
%       IP address} \\
%     \multicolumn{8}{@{\vrule}c@{\vrule}}{} \\
%     \hline
%     \multicolumn{8}{@{\vrule}c@{\vrule}}{} \\
%     \multicolumn{8}{@{\vrule}c@{\vrule}}{
%       destination} \\
%     \multicolumn{8}{@{\vrule}c@{\vrule}}{
%      IP address} \\
%     \multicolumn{8}{@{\vrule}c@{\vrule}}{} \\
%     \hline
%     \end{tabular}
%     \caption{L'intestazione o \textit{header} di IPv6}
%     \label{tab:IP_ipv6head}
%   \end{center}
% \end{table}

\begin{figure}[!htb]
  \centering
  \includegraphics[width=10cm]{img/ipv6_head}
  \caption{L'intestazione o \textit{header} di IPv6.}
  \label{fig:IP_ipv6head}
\end{figure}


Come si può notare l'intestazione di IPv6 diventa di dimensione fissa, pari a
40 byte, contro una dimensione (minima, in assenza di opzioni) di 20 byte per
IPv4; un semplice raddoppio nonostante lo spazio destinato agli indirizzi sia
quadruplicato, questo grazie a una notevole semplificazione che ha ridotto il
numero dei campi da 12 a 8.

\begin{table}[htb]
  \begin{center}
  \footnotesize
    \begin{tabular}{|l|c|p{8cm}|}
      \hline
      \textbf{Nome} & \textbf{Bit} & \textbf{Significato} \\
      \hline
      \hline
      \textit{version}       & 4& La \textsl{versione}, nel caso specifico vale
                                  sempre 6.\\ 
      \textit{priority}      & 4& La \textsl{priorità}, vedi
                                  sez.~\ref{sec:IPv6_prio}.\\
      \textit{flow label}    &24& L'\textsl{etichetta di flusso}, vedi
                                  sez.~\ref{sec:IP_ipv6_flow}.\\ 
      \textit{payload length}&16& La \textsl{lunghezza del carico}, cioè del
                                  corpo dei dati che segue l'intestazione, in
                                  byte. \\ 
      \textit{next header}   & 8& L'\textsl{intestazione successiva}, 
                                  identifica il tipo di pacchetto che segue
                                  l'intestazione di IPv6, ed usa gli stessi
                                  valori del campo protocollo
                                  nell'intestazione di IPv4.\\ 
      \textit{hop limit}     & 8& Il \textsl{limite di salti}, ha lo stesso 
                                  significato del \textit{time to live} 
                                  nell'intestazione di IPv4.\\ 
      \textit{source IP}     &128&L'\textsl{indirizzo di origine}.\\
      \textit{destination IP}&128&L'\textsl{indirizzo di destinazione}.\\
      \hline
    \end{tabular}
    \caption{Legenda per il significato dei campi dell'intestazione di IPv6}
    \label{tab:IP_ipv6field}
  \end{center}
\end{table}

Abbiamo già anticipato in sez.~\ref{sec:IP_ipv6over} uno dei criteri
principali nella progettazione di IPv6 è stato quello di ridurre al minimo il
tempo di elaborazione dei pacchetti da parte dei router, un confronto con
l'intestazione di IPv4 (vedi fig.~\ref{fig:IP_ipv4_head}) mostra le seguenti
differenze:

\begin{itemize}
\item è stato eliminato il campo \textit{header length} in quanto le opzioni
  sono state tolte dall'intestazione che ha così dimensione fissa; ci possono
  essere più intestazioni opzionali (\textsl{intestazioni di estensione}, vedi
  sez.~\ref{sec:IP_ipv6_extens}), ciascuna delle quali avrà un suo campo di
  lunghezza all'interno.
\item l'intestazione e gli indirizzi sono allineati a 64 bit, questo rende più
  veloce il processo da parte di computer con processori a 64 bit.
\item i campi per gestire la frammentazione (\textit{identification},
  \textit{flag} e \textit{fragment offset}) sono stati eliminati; questo
  perché la frammentazione è un'eccezione che non deve rallentare
  l'elaborazione dei pacchetti nel caso normale.
\item è stato eliminato il campo \textit{checksum} in quanto tutti i
  protocolli di livello superiore (TCP, UDP e ICMPv6) hanno un campo di
  checksum che include, oltre alla loro intestazione e ai dati, pure i campi
  \textit{payload length}, \textit{next header}, e gli indirizzi di origine e
  di destinazione; una checksum esiste anche per la gran parte protocolli di
  livello inferiore (anche se quelli che non lo hanno, come SLIP, non possono
  essere usati con grande affidabilità); con questa scelta si è ridotto di
  molto il tempo di elaborazione dato che i router non hanno più la necessità
  di ricalcolare la checksum ad ogni passaggio di un pacchetto per il
  cambiamento del campo \textit{hop limit}.
\item è stato eliminato il campo \textit{type of service}, che praticamente
  non è mai stato utilizzato; una parte delle funzionalità ad esso delegate
  sono state reimplementate (vedi il campo \textit{priority} al prossimo
  punto) con altri metodi.
\item è stato introdotto un nuovo campo \textit{flow label}, che viene usato,
  insieme al campo \textit{priority} (che recupera i bit di precedenza del
  campo \textit{type of service}) per implementare la gestione di una
  ``\textsl{qualità di servizio}'' (vedi sez.~\ref{sec:IP_ipv6_qos}) che
  permette di identificare i pacchetti appartenenti a un ``\textsl{flusso}''
  di dati per i quali si può provvedere un trattamento speciale.
\end{itemize}

Oltre alle differenze precedenti, relative ai singoli campi nell'intestazione,
ulteriori caratteristiche che diversificano il comportamento di IPv4 da
quello di IPv6 sono le seguenti:

\begin{itemize}
\item il \textit{broadcasting} non è previsto in IPv6, le applicazioni che lo
  usano dovono essere reimplementate usando il \textit{multicasting} (vedi
  sez.~\ref{sec:IP_ipv6_multicast}), che da opzionale diventa obbligatorio.
\item è stato introdotto un nuovo tipo di indirizzi, gli \textit{anycast}.
\item i router non possono più frammentare i pacchetti lungo il cammino, la
  frammentazione di pacchetti troppo grandi potrà essere gestita solo ai
  capi della comunicazione (usando un'apposita estensione vedi
  sez.~\ref{sec:IP_ipv6_extens}).
\item IPv6 richiede il supporto per il \textit{path MTU discovery} (cioè il
  protocollo per la selezione della massima lunghezza del pacchetto); seppure
  questo sia in teoria opzionale, senza di esso non sarà possibile inviare
  pacchetti più larghi della dimensione minima (576 byte).
\end{itemize}

\subsection{Gli indirizzi di IPv6}
\label{sec:IP_ipv6_addr}

Come già abbondantemente anticipato la principale novità di IPv6 è
costituita dall'ampliamento dello spazio degli indirizzi, che consente di avere
indirizzi disponibili in un numero dell'ordine di quello degli atomi che
costituiscono la terra. 

In realtà l'allocazione di questi indirizzi deve tenere conto della
necessità di costruire delle gerarchie che consentano un instradamento
rapido ed efficiente dei pacchetti, e flessibilità nel dispiegamento delle
reti, il che comporta una riduzione drastica dei numeri utilizzabili; uno
studio sull'efficienza dei vari sistemi di allocazione usati in altre
architetture (come i sistemi telefonici) è comunque giunto alla conclusione
che anche nella peggiore delle ipotesi IPv6 dovrebbe essere in grado di
fornire più di un migliaio di indirizzi per ogni metro quadro della
superficie terrestre.


\subsection{La notazione}
\label{sec:IP_ipv6_notation}
Con un numero di bit quadruplicato non è più possibile usare la notazione
coi numeri decimali di IPv4 per rappresentare un numero IP. Per questo gli
indirizzi di IPv6 sono in genere scritti come sequenze di otto numeri
esadecimali di 4 cifre (cioè a gruppi di 16 bit) usando i due punti come
separatore; cioè qualcosa del tipo
\texttt{1080:0000:0000:0008:0800:ba98:2078:e3e3}.

Visto che la notazione resta comunque piuttosto pesante esistono alcune
abbreviazioni: si può evitare di scrivere gli zeri iniziali delle singole
cifre, abbreviando l'indirizzo precedente in
\texttt{1080:0:0:8:800:ba98:2078:e3e3}; se poi un intero è zero si può
omettere del tutto, così come un insieme di zeri (ma questo solo una volta per
non generare ambiguità) per cui il precedente indirizzo si può scrivere anche
come \texttt{1080::8:800:ba98:2078:e3e3}.

Infine per scrivere un indirizzo IPv4 all'interno di un indirizzo IPv6 si
può usare la vecchia notazione con i punti, per esempio
\texttt{::192.84.145.138}.

\begin{table}[htb]
  \centering 
  \footnotesize
  \begin{tabular}{|l|l|l|}
    \hline
    \centering \textbf{Tipo di indirizzo}
    & \centering \textbf{Prefisso} & {\centering \textbf{Frazione}} \\
    \hline
    \hline
    riservato & \texttt{0000 0000} & 1/256 \\
    non assegnato  & \texttt{0000 0001} & 1/256 \\
    \hline
    riservato per NSAP & \texttt{0000 001} & 1/128\\
    riservato per IPX & \texttt{0000 010} & 1/128\\
    \hline
    non assegnato  & \texttt{0000 011} & 1/128 \\
    non assegnato  & \texttt{0000 1} & 1/32 \\
    non assegnato  & \texttt{0001} & 1/16 \\
    \hline
    provider-based & \texttt{001} & 1/8\\
    \hline
    non assegnato  & \texttt{010} & 1/8 \\
    non assegnato  & \texttt{011} & 1/8 \\
    geografic-based& \texttt{100} & 1/8 \\
    non assegnato  & \texttt{101} & 1/8 \\
    non assegnato  & \texttt{110} & 1/8 \\
    non assegnato  & \texttt{1110} & 1/16 \\
    non assegnato  & \texttt{1111 0} & 1/32 \\
    non assegnato  & \texttt{1111 10} & 1/64 \\
    non assegnato  & \texttt{1111 110} & 1/128 \\
    non assegnato  & \texttt{1111 1110 0} & 1/512 \\
    \hline
    unicast link-local & \texttt{1111 1110 10} & 1/1024 \\
    unicast site-local & \texttt{1111 1110 11} & 1/1024 \\
    \hline
    \hline
    \textit{multicast} & \texttt{1111 1111} & 1/256 \\
    \hline
  \end{tabular}
  \caption{Classificazione degli indirizzi IPv6 a seconda dei bit più 
    significativi}
  \label{tab:IP_ipv6addr}
\end{table}


\subsection{La architettura degli indirizzi di IPv6}
\label{sec:IP_ipv6_addr_arch}

Come per IPv4 gli indirizzi sono identificatori per una singola (indirizzi
\textit{unicast}) o per un insieme (indirizzi \textit{multicast} e
\textit{anycast}) di interfacce di rete.

Gli indirizzi sono sempre assegnati all'interfaccia, non al nodo che la
ospita; dato che ogni interfaccia appartiene ad un nodo quest'ultimo può
essere identificato attraverso uno qualunque degli indirizzi unicast delle sue
interfacce. A una interfaccia possono essere associati anche più indirizzi.

IPv6 presenta tre tipi diversi di indirizzi: due di questi, gli indirizzi
\textit{unicast} e \textit{multicast} hanno le stesse caratteristiche che in
IPv4, un terzo tipo, gli indirizzi \textit{anycast} è completamente nuovo.  In
IPv6 non esistono più gli indirizzi \textit{broadcast}, la funzione di questi
ultimi deve essere reimplementata con gli indirizzi \textit{multicast}.

Gli indirizzi \textit{unicast} identificano una singola interfaccia: i
pacchetti mandati ad un tale indirizzo verranno inviati a quella interfaccia,
gli indirizzi \textit{anycast} identificano un gruppo di interfacce tale che
un pacchetto mandato a uno di questi indirizzi viene inviato alla più vicina
(nel senso di distanza di routing) delle interfacce del gruppo, gli indirizzi
\textit{multicast} identificano un gruppo di interfacce tale che un pacchetto
mandato a uno di questi indirizzi viene inviato a tutte le interfacce del
gruppo.

In IPv6 non ci sono più le classi ma i bit più significativi indicano il tipo
di indirizzo; in tab.~\ref{tab:IP_ipv6addr} sono riportati i valori di detti
bit e il tipo di indirizzo che loro corrispondente.  I bit più significativi
costituiscono quello che viene chiamato il \textit{format prefix} ed è sulla
base di questo che i vari tipi di indirizzi vengono identificati.  Come si
vede questa architettura di allocazione supporta l'allocazione di indirizzi
per i provider, per uso locale e per il \textit{multicast}; inoltre è stato
riservato lo spazio per indirizzi NSAP, IPX e per le connessioni; gran parte
dello spazio (più del 70\%) è riservato per usi futuri.

Si noti infine che gli indirizzi \textit{anycast} non sono riportati in
tab.~\ref{tab:IP_ipv6addr} in quanto allocati al di fuori dello spazio di
allocazione degli indirizzi unicast.

\subsection{Indirizzi unicast \textit{provider-based}}
\label{sec:IP_ipv6_unicast}

Gli indirizzi \textit{provider-based} sono gli indirizzi usati per le
comunicazioni globali, questi sono definiti
nell'\href{http://www.ietf.org/rfc/rfc2073.txt}{RFC~2073} e sono gli
equivalenti degli attuali indirizzi delle classi da A a C.

L'autorità che presiede all'allocazione di questi indirizzi è la IANA; per
evitare i problemi di crescita delle tabelle di instradamento e una procedura
efficiente di allocazione la struttura di questi indirizzi è organizzata fin
dall'inizio in maniera gerarchica; pertanto lo spazio di questi indirizzi è
stato suddiviso in una serie di campi secondo lo schema riportato in
tab.~\ref{tab:IP_ipv6_unicast}.

\begin{table}[htb]
  \centering
  \footnotesize
  \begin{tabular} {@{\vrule}p{6mm}
      @{\vrule}p{16mm}@{\vrule}p{24mm}
      @{\vrule}p{30mm}@{\vrule}c@{\vrule}}
    \multicolumn{1}{@{}c@{}}{3}&\multicolumn{1}{c}{5 bit}&
    \multicolumn{1}{c}{$n$ bit}&\multicolumn{1}{c}{$56-n$ bit}&
    \multicolumn{1}{c}{64 bit} \\
    \hline
    \omit\vrule\hfill\vrule&\hspace{16mm} & & &\omit\hspace{76mm}\hfill\vrule\\
    \centering 010&
    \centering \textsl{Registry Id}&
    \centering \textsl{Provider Id}& 
    \centering \textsl{Subscriber Id}& 
    \textsl{Intra-Subscriber} \\
    \omit\vrule\hfill\vrule& & & &\omit\hspace{6mm}\hfill\vrule\\ 
    \hline
  \end{tabular}
\caption{Formato di un indirizzo unicast \textit{provider-based}.}
\label{tab:IP_ipv6_unicast}
\end{table}

Al livello più alto la IANA può delegare l'allocazione a delle autorità
regionali (i Regional Register) assegnando ad esse dei blocchi di indirizzi; a
queste autorità regionali è assegnato un Registry Id che deve seguire
immediatamente il prefisso di formato. Al momento sono definite tre registri
regionali (INTERNIC, RIPE NCC e APNIC), inoltre la IANA si è riservata la
possibilità di allocare indirizzi su base regionale; pertanto sono previsti
i seguenti possibili valori per il \textsl{Registry Id};
gli altri valori restano riservati per la IANA.
\begin{table}[htb]
  \centering 
  \footnotesize
    \begin{tabular}{|l|l|l|}
      \hline
      \textbf{Regione} & \textbf{Registro} & \textbf{Id} \\
      \hline
      \hline
      Nord America &INTERNIC & \texttt{11000} \\
      Europa & RIPE NCC & \texttt{01000} \\
      Asia & APNIC & \texttt{00100} \\
      Multi-regionale & IANA &\texttt{10000} \\
      \hline
    \end{tabular}
    \caption{Valori dell'identificativo dei 
      Regional Register allocati ad oggi.}
    \label{tab:IP_ipv6_regid}
\end{table}

L'organizzazione degli indirizzi prevede poi che i due livelli successivi, di
suddivisione fra \textit{Provider Id}, che identifica i grandi fornitori di
servizi, e \textit{Subscriber Id}, che identifica i fruitori, sia gestita dai
singoli registri regionali. Questi ultimi dovranno definire come dividere lo
spazio di indirizzi assegnato a questi due campi (che ammonta a un totale di
56~bit), definendo lo spazio da assegnare al \textit{Provider Id} e
al \textit{Subscriber Id}, ad essi spetterà inoltre anche l'allocazione dei
numeri di \textit{Provider Id} ai singoli fornitori, ai quali sarà delegata
l'autorità di allocare i \textit{Subscriber Id} al loro interno.

L'ultimo livello è quello \textit{Intra-subscriber} che è lasciato alla
gestione dei singoli fruitori finali, gli indirizzi \textit{provider-based}
lasciano normalmente gli ultimi 64~bit a disposizione per questo livello, la
modalità più immediata è quella di usare uno schema del tipo mostrato in
tab.~\ref{tab:IP_ipv6_uninterf} dove l'\textit{Interface Id} è dato dal
MAC-address a 48~bit dello standard Ethernet, scritto in genere nell'hardware
delle scheda di rete, e si usano i restanti 16~bit per indicare la sottorete.

\begin{table}[htb]
  \centering
  \footnotesize
  \begin{tabular} {@{\vrule}p{64mm}@{\vrule}p{16mm}@{\vrule}c@{\vrule}}
    \multicolumn{1}{c}{64 bit}&\multicolumn{1}{c}{16 bit}&
    \multicolumn{1}{c}{48 bit}\\
    \hline
    \omit\vrule\hfill\vrule&\hspace{16mm}&\omit\hspace{48mm}\hfill\vrule\\ 
    \centering \textsl{Subscriber Prefix}& 
    \centering \textsl{Subnet Id}&
    \textsl{Interface Id}\\
    \omit\vrule\hfill\vrule& &\omit\hspace{6mm}\hfill\vrule\\ 
    \hline
  \end{tabular}
\caption{Formato del campo \textit{Intra-subscriber} per un indirizzo unicast
  \textit{provider-based}.}
\label{tab:IP_ipv6_uninterf}
\end{table}

Qualora si dovesse avere a che fare con una necessità di un numero più
elevato di sotto-reti, il precedente schema andrebbe modificato, per evitare
l'enorme spreco dovuto all'uso dei MAC-address, a questo scopo si possono
usare le capacità di auto-configurazione di IPv6 per assegnare indirizzi
generici con ulteriori gerarchie per sfruttare efficacemente tutto lo spazio
di indirizzi.

Un registro regionale può introdurre un ulteriore livello nella gerarchia
degli indirizzi, allocando dei blocchi per i quali delegare l'autorità a dei
registri nazionali, quest'ultimi poi avranno il compito di gestire la
attribuzione degli indirizzi per i fornitori di servizi nell'ambito del/i
paese coperto dal registro nazionale con le modalità viste in precedenza.
Una tale ripartizione andrà effettuata all'interno dei soliti 56~bit come
mostrato in tab.~\ref{tab:IP_ipv6_uninaz}.

\begin{table}[htb]
  \centering
  \footnotesize
  \begin{tabular} {@{\vrule}p{3mm}
      @{\vrule}p{10mm}@{\vrule}p{12mm}@{\vrule}p{18mm}
      @{\vrule}p{18mm}@{\vrule}c@{\vrule}}
    \multicolumn{1}{@{}c@{}}{3}&\multicolumn{1}{c}{5 bit}&
    \multicolumn{1}{c}{n bit}&\multicolumn{1}{c}{m bit}&
    \multicolumn{1}{c}{56-n-m bit}&\multicolumn{1}{c}{64 bit} \\
    \hline
    \omit\vrule\hfill\vrule& & & & &\omit\hspace{64mm}\hfill\vrule\\
    \centering \texttt{3}&
    \centering \textsl{Reg.}&
    \centering \textsl{Naz.}&
    \centering \textsl{Prov.}& 
    \centering \textsl{Subscr.}& 
    \textsl{Intra-Subscriber} \\
    \omit\vrule\hfill\vrule &&&&&\\ 
    \hline
  \end{tabular}
\caption{Formato di un indirizzo unicast \textit{provider-based} che prevede
      un registro nazionale.}
\label{tab:IP_ipv6_uninaz}
\end{table}


\subsection{Indirizzi ad uso locale}
\label{sec:IP_ipv6_linksite}

Gli indirizzi ad uso locale sono indirizzi unicast che sono instradabili solo
localmente (all'interno di un sito o di una sottorete), e possono avere una
unicità locale o globale.

Questi indirizzi sono pensati per l'uso all'interno di un sito per mettere su
una comunicazione locale immediata, o durante le fasi di auto-configurazione
prima di avere un indirizzo globale.

\begin{table}[htb]
  \centering
  \footnotesize
  \begin{tabular} {@{\vrule}p{10mm}@{\vrule}p{54mm}@{\vrule}c@{\vrule}}
    \multicolumn{1}{c}{10} &\multicolumn{1}{c}{54 bit} & 
    \multicolumn{1}{c}{64 bit} \\
    \hline
    \omit\vrule\hfill\vrule & & \omit\hspace{64mm}\hfill\vrule\\
    \centering \texttt{FE80}& 
    \centering\texttt{0000 .   .   .   .   . 0000} &
    Interface Id \\
    \omit\vrule\hfill\vrule & &\\
    \hline
\end{tabular}
\caption{Formato di un indirizzo \textit{link-local}.}
\label{tab:IP_ipv6_linklocal}
\end{table}

Ci sono due tipi di indirizzi, \textit{link-local} e \textit{site-local}. Il
primo è usato per un singolo link; la struttura è mostrata in
tab.~\ref{tab:IP_ipv6_linklocal}, questi indirizzi iniziano sempre con un
valore nell'intervallo \texttt{FE80}--\texttt{FEBF} e vengono in genere usati
per la configurazione automatica dell'indirizzo al bootstrap e per la ricerca
dei vicini (vedi \ref{sec:IP_ipv6_autoconf}); un pacchetto che abbia tale
indirizzo come sorgente o destinazione non deve venire ritrasmesso dai router,
sono gli indirizzi che identificano la macchina sulla rete locale, per questo
sono chiamati in questo modo, in quanto sono usati solo su di essa.

Un indirizzo \textit{site-local} invece è usato per l'indirizzamento
all'interno di un sito che non necessita di un prefisso globale; la struttura
è mostrata in tab.~\ref{tab:IP_ipv6_sitelocal}, questi indirizzi iniziano
sempre con un valore nell'intervallo \texttt{FEC0}--\texttt{FEFF} e non devono
venire ritrasmessi dai router all'esterno del sito stesso; sono in sostanza
gli equivalenti degli indirizzi riservati per reti private definiti su IPv4.
Per entrambi gli indirizzi il campo \textit{Interface Id} è un identificatore
che deve essere unico nel dominio in cui viene usato, un modo immediato per
costruirlo è quello di usare il MAC-address delle schede di rete.
 
\begin{table}[!h]
  \centering
  \footnotesize
  \begin{tabular} {@{\vrule}p{10mm}@{\vrule}p{38mm}@{\vrule}p{16mm}
      @{\vrule}c@{\vrule}}
    \multicolumn{1}{c}{10} &\multicolumn{1}{c}{38 bit} & 
    \multicolumn{1}{c}{16 bit} &\multicolumn{1}{c}{64 bit} \\
    \hline
    \omit\vrule\hfill\vrule& & & \omit\hspace{64mm}\hfill\vrule\\
    \centering \texttt{FEC0}& 
    \centering \texttt{0000 .   .   . 0000}& 
    \centering Subnet Id &
    Interface Id\\
    \omit\vrule\hfill\vrule& & &\\
    \hline
\end{tabular}
\caption{Formato di un indirizzo \textit{site-local}.}
\label{tab:IP_ipv6_sitelocal}
\end{table}

Gli indirizzi di uso locale consentono ad una organizzazione che non è
(ancora) connessa ad Internet di operare senza richiedere un prefisso globale,
una volta che in seguito l'organizzazione venisse connessa a Internet
potrebbe continuare a usare la stessa suddivisione effettuata con gli
indirizzi \textit{site-local} utilizzando un prefisso globale e la
rinumerazione degli indirizzi delle singole macchine sarebbe automatica.

\subsection{Indirizzi riservati}
\label{sec:IP_ipv6_reserved}

Alcuni indirizzi sono riservati per scopi speciali, in particolare per scopi
di compatibilità.

Un primo tipo sono gli indirizzi \textit{IPv4 mappati su IPv6} (mostrati in
tab.~\ref{tab:IP_ipv6_map}), questo sono indirizzi unicast che vengono usati
per consentire ad applicazioni IPv6 di comunicare con host capaci solo di
IPv4; questi sono ad esempio gli indirizzi generati da un DNS quando l'host
richiesto supporta solo IPv4; l'uso di un tale indirizzo in un socket IPv6
comporta la generazione di un pacchetto IPv4 (ovviamente occorre che sia IPv4
che IPv6 siano supportati sull'host di origine).

\begin{table}[!htb]
  \centering
  \footnotesize
  \begin{tabular} {@{\vrule}p{80mm}@{\vrule}p{16mm}@{\vrule}c@{\vrule}}
    \multicolumn{1}{c}{80 bit} &\multicolumn{1}{c}{16 bit} & 
    \multicolumn{1}{c}{32 bit} \\
    \hline
    \omit\vrule\hfill\vrule& &\omit\hspace{32mm}\hfill\vrule\\ 
    \centering
    \texttt{0000 .   .   .   .   .   .   .   .   .   .   .   . 0000} & 
    \centering\texttt{FFFF} &
    IPv4 address \\
    \omit\vrule\hfill\vrule& &\\ 
    \hline
\end{tabular}
\caption{Formato di un indirizzo IPV4 mappato su IPv6.}
\label{tab:IP_ipv6_map}
\end{table}

Un secondo tipo di indirizzi di compatibilità sono gli \textsl{IPv4
  compatibili IPv6} (vedi tab.~\ref{tab:IP_ipv6_comp}) usati nella transizione
da IPv4 a IPv6: quando un nodo che supporta sia IPv6 che IPv4 non ha un router
IPv6 deve usare nel DNS un indirizzo di questo tipo, ogni pacchetto IPv6
inviato a un tale indirizzo verrà automaticamente incapsulato in IPv4.

\begin{table}[htb]
  \centering
  \footnotesize
  \begin{tabular} {@{\vrule}p{80mm}@{\vrule}p{16mm}@{\vrule}p{32mm}@{\vrule}}
    \multicolumn{1}{c}{80 bit} &\multicolumn{1}{c}{16 bit} & 
    \multicolumn{1}{c}{32 bit} \\
    \hline
    \omit\vrule\hfill\vrule& &\omit\hspace{32mm}\hfill\vrule\\ 
    \centering
    \texttt{0000 .   .   .   .   .   .   .   .   .   .   .   . 0000} & 
    \centering\texttt{0000} &
    \parbox{32mm}{\centering IPv4 address} \\
    \omit\vrule\hfill\vrule& &\\ 
    \hline
\end{tabular}
\caption{Formato di un indirizzo IPV4 mappato su IPv6.}
\label{tab:IP_ipv6_comp}
\end{table}

Altri indirizzi speciali sono il \textit{loopback address}, costituito da 127
zeri ed un uno (cioè \texttt{::1}) e l'\textsl{indirizzo generico}
costituito da tutti zeri (scritto come \texttt{0::0} o ancora più
semplicemente come \texttt{:}) usato in genere quando si vuole indicare
l'accettazione di una connessione da qualunque host.

\subsection{Multicasting}
\label{sec:IP_ipv6_multicast}

\itindbeg{multicast}

Gli indirizzi \textit{multicast} sono usati per inviare un pacchetto a un
gruppo di interfacce; l'indirizzo identifica uno specifico gruppo di
\textit{multicast} e il pacchetto viene inviato a tutte le interfacce di detto
gruppo.  Un'interfaccia può appartenere ad un numero qualunque numero di
gruppi di \textit{multicast}. Il formato degli indirizzi \textit{multicast} è
riportato in tab.~\ref{tab:IP_ipv6_multicast}:

\begin{table}[htb]
  \centering
  \footnotesize
  \begin{tabular} {@{\vrule}p{12mm}
      @{\vrule}p{6mm}@{\vrule}p{6mm}@{\vrule}c@{\vrule}}
    \multicolumn{1}{c}{8}&\multicolumn{1}{c}{4}&
    \multicolumn{1}{c}{4}&\multicolumn{1}{c}{112 bit} \\
    \hline
    \omit\vrule\hfill\vrule& & & \omit\hspace{104mm}\hfill\vrule\\
    \centering\texttt{FF}& 
    \centering flag &
    \centering scop& 
    Group Id\\
    \omit\vrule\hfill\vrule &&&\\ 
    \hline
  \end{tabular}
\caption{Formato di un indirizzo \textit{multicast}.}
\label{tab:IP_ipv6_multicast}
\end{table}

Il prefisso di formato per tutti gli indirizzi \textit{multicast} è
\texttt{FF}, ad esso seguono i due campi il cui significato è il seguente:

\begin{itemize}
\item \textsl{flag}: un insieme di 4 bit, di cui i primi tre sono riservati e
  posti a zero, l'ultimo è zero se l'indirizzo è permanente (cioè un
  indirizzo noto, assegnato dalla IANA), ed è uno se invece l'indirizzo è
  transitorio.
\item \textsl{scop} è un numero di quattro bit che indica il raggio di
  validità dell'indirizzo, i valori assegnati per ora sono riportati in
  tab.~\ref{tab:IP_ipv6_multiscope}.
\end{itemize}



\begin{table}[!htb]
  \centering 
  \footnotesize
  \begin{tabular}[c]{|c|l|c|l|}
    \hline
    \multicolumn{4}{|c|}{\bf Gruppi di multicast} \\
    \hline
    \hline
    0 & riservato & 8 & organizzazione locale \\
    1 & nodo locale & 9 & non assegnato \\
    2 & collegamento locale & A & non assegnato \\
    3 & non assegnato & B & non assegnato \\
    4 & non assegnato & C & non assegnato \\ 
    5 & sito locale & D & non assegnato \\
    6 & non assegnato & E & globale \\
    7 & non assegnato & F & riservato \\
    \hline
  \end{tabular}
\caption{Possibili valori del campo \textsl{scop} di un indirizzo
  \textit{multicast}.} 
\label{tab:IP_ipv6_multiscope}
\end{table}

Infine l'ultimo campo identifica il gruppo di \textit{multicast}, sia
permanente che transitorio, all'interno del raggio di validità del medesimo.
Alcuni indirizzi \textit{multicast}, riportati in tab.~\ref{tab:multiadd} sono
già riservati per il funzionamento della rete.

\begin{table}[!htb]
  \centering 
  \footnotesize
  \begin{tabular}[c]{|l|l|l|}
    \hline
    \textbf{Uso}& \textbf{Indirizzi riservati} & \textbf{Definizione}\\
    \hline 
    \hline 
    all-nodes       & \texttt{FFxx:0:0:0:0:0:0:1}  & 
                      \href{http://www.ietf.org/rfc/rfc1970.txt}{RFC~1970} \\
    all-routers     & \texttt{FFxx:0:0:0:0:0:0:2}  & 
                      \href{http://www.ietf.org/rfc/rfc1970.txt}{RFC~1970} \\
    all-rip-routers & \texttt{FFxx:0:0:0:0:0:0:9}  & 
                      \href{http://www.ietf.org/rfc/rfc2080.txt}{RFC~2080} \\
    all-cbt-routers & \texttt{FFxx:0:0:0:0:0:0:10} & \\
    reserved        & \texttt{FFxx:0:0:0:0:0:1:0}  & IANA \\
    link-name       & \texttt{FFxx:0:0:0:0:0:1:1}  &  \\
    all-dhcp-agents & \texttt{FFxx:0:0:0:0:0:1:2}  & \\
    all-dhcp-servers& \texttt{FFxx:0:0:0:0:0:1:3}  & \\
    all-dhcp-relays & \texttt{FFxx:0:0:0:0:0:1:4}  & \\
    solicited-nodes & \texttt{FFxx:0:0:0:0:1:0:0}  & 
                      \href{http://www.ietf.org/rfc/rfc1970.txt}{RFC~1970} \\
    \hline
  \end{tabular}
\caption{Gruppi di \textit{multicast} predefiniti.}
\label{tab:multiadd}
\end{table}

L'utilizzo del campo di \textit{scope} e di questi indirizzi predefiniti serve
a recuperare le funzionalità del \textit{broadcasting} (ad esempio inviando un
pacchetto all'indirizzo \texttt{FF02:0:0:0:0:0:0:1} si raggiungono tutti i
nodi locali).

\itindend{multicast}

\subsection{Indirizzi \textit{anycast}}
\label{sec:IP_anycast}

Gli indirizzi \textit{anycast} sono indirizzi che vengono assegnati ad un
gruppo di interfacce: un pacchetto indirizzato a questo tipo di indirizzo
viene inviato al componente del gruppo più ``\textsl{vicino}'' secondo la
distanza di instradamento calcolata dai router.

Questi indirizzi sono allocati nello stesso spazio degli indirizzi unicast,
usando uno dei formati disponibili, e per questo, sono da essi assolutamente
indistinguibili. Quando un indirizzo unicast viene assegnato a più interfacce
(trasformandolo in un anycast) il computer su cui è l'interfaccia deve essere
configurato per tener conto del fatto.

Gli indirizzi anycast consentono a un nodo sorgente di inviare pacchetti a una
destinazione su un gruppo di possibili interfacce selezionate. La sorgente non
deve curarsi di come scegliere l'interfaccia più vicina, compito che tocca al
sistema di instradamento (in sostanza la sorgente non ha nessun controllo
sulla selezione).

Gli indirizzi anycast, quando vengono usati come parte di una sequenza di
instradamento, consentono ad esempio ad un nodo di scegliere quale fornitore
vuole usare (configurando gli indirizzi anycast per identificare i router di
uno stesso provider).

Questi indirizzi pertanto possono essere usati come indirizzi intermedi in una
intestazione di instradamento o per identificare insiemi di router connessi a
una particolare sottorete, o che forniscono l'accesso a un certo sotto
dominio.

L'idea alla base degli indirizzi anycast è perciò quella di utilizzarli per
poter raggiungere il fornitore di servizio più vicino; ma restano aperte tutta
una serie di problematiche, visto che una connessione con uno di questi
indirizzi non è possibile, dato che per una variazione delle distanze di
routing non è detto che due pacchetti successivi finiscano alla stessa
interfaccia.

La materia è pertanto ancora controversa e in via di definizione.


\subsection{Le estensioni}
\label{sec:IP_ipv6_extens}

Come già detto in precedenza IPv6 ha completamente cambiato il trattamento
delle opzioni; queste ultime infatti sono state tolte dall'intestazione del
pacchetto, e poste in apposite \textsl{intestazioni di estensione} (o
\textit{extension header}) poste fra l'intestazione di IPv6 e l'intestazione
del protocollo di trasporto.

Per aumentare la velocità di elaborazione, sia dei dati del livello seguente
che di ulteriori opzioni, ciascuna estensione deve avere una lunghezza
multipla di 8 byte per mantenere l'allineamento a 64~bit di tutti le
intestazioni seguenti.

Dato che la maggior parte di queste estensioni non sono esaminate dai router
durante l'instradamento e la trasmissione dei pacchetti, ma solo all'arrivo
alla destinazione finale, questa scelta ha consentito un miglioramento delle
prestazioni rispetto a IPv4 dove la presenza di un'opzione comportava l'esame
di tutte quante.

Un secondo miglioramento è che rispetto a IPv4 le opzioni possono essere di
lunghezza arbitraria e non limitate a 40 byte; questo, insieme al modo in cui
vengono trattate, consente di utilizzarle per scopi come l'autenticazione e la
sicurezza, improponibili con IPv4.

Le estensioni definite al momento sono le seguenti:
\begin{itemize}
\item \textbf{Hop by hop} devono seguire immediatamente l'intestazione
  principale; indicano le opzioni che devono venire processate ad ogni
  passaggio da un router, fra di esse è da menzionare la \textit{jumbo
    payload} che segnala la presenza di un pacchetto di dati di dimensione
  superiore a 65535 byte.
\item \textbf{Destination options} opzioni che devono venire esaminate al nodo
  di ricevimento, nessuna di esse è tuttora definita.
\item \textbf{Routing} definisce una \textit{source route} (come la analoga
  opzione di IPv4) cioè una lista di indirizzi IP di nodi per i quali il
  pacchetto deve passare. 
\item \textbf{Fragmentation} viene generato automaticamente quando un host
  vuole frammentare un pacchetto, ed è riprocessato automaticamente alla
  destinazione che riassembla i frammenti.
\item \textbf{Authentication} gestisce l'autenticazione e il controllo di
  integrità dei pacchetti; è documentato
  dall'\href{http://www.ietf.org/rfc/rfc1826.txt}{RFC~1826}.
\item \textbf{Encapsulation} serve a gestire la segretezza del contenuto
  trasmesso; è documentato
  dall'\href{http://www.ietf.org/rfc/rfc1827.txt}{RFC~1827}.
\end{itemize}

La presenza di opzioni è rilevata dal valore del campo \textit{next header}
che indica qual è l'intestazione successiva a quella di IPv6; in assenza di
opzioni questa sarà l'intestazione di un protocollo di trasporto del livello
superiore, per cui il campo assumerà lo stesso valore del campo
\textit{protocol} di IPv4, altrimenti assumerà il valore dell'opzione
presente; i valori possibili sono riportati in
tab.~\ref{tab:IP_ipv6_nexthead}.

\begin{table}[htb]
  \begin{center}
    \footnotesize
    \begin{tabular}{|c|l|l|}
      \hline
      \textbf{Valore} & \textbf{Keyword} & \textbf{Tipo di protocollo} \\
      \hline
      \hline
      0  &      & Riservato.\\
         & HBH  & Hop by Hop.\\
      1  & ICMP & Internet Control Message (IPv4 o IPv6).\\
      2  & IGMP & Internet Group Management (IPv4).\\
      3  & GGP  & Gateway-to-Gateway.\\
      4  & IP   & IP in IP (IPv4 encapsulation).\\
      5  & ST   & Stream.\\
      6  & TCP  & Trasmission Control.\\
      17 & UDP  & User Datagram.\\
      43 & RH   & Routing Header (IPv6).\\
      44 & FH   & Fragment Header (IPv6).\\
      45 & IDRP & Inter Domain Routing.\\
      51 & AH   & Authentication Header (IPv6).\\
      52 & ESP  & Encrypted Security Payload (IPv6).\\
      59 & Null & No next header (IPv6).\\
      88 & IGRP & Internet Group Routing.\\
      89 & OSPF & Open Short Path First.\\
      255&      & Riservato.\\
    \hline
    \end{tabular}
    \caption{Tipi di protocolli e intestazioni di estensione}
    \label{tab:IP_ipv6_nexthead}
  \end{center}
\end{table}

Questo meccanismo permette la presenza di più opzioni in successione prima
del pacchetto del protocollo di trasporto; l'ordine raccomandato per le
estensioni è quello riportato nell'elenco precedente con la sola differenza
che le opzioni di destinazione sono inserite nella posizione ivi indicata solo
se, come per il tunnelling, devono essere esaminate dai router, quelle che
devono essere esaminate solo alla destinazione finale vanno in coda.


\subsection{Qualità di servizio}
\label{sec:IP_ipv6_qos}

Una delle caratteristiche innovative di IPv6 è quella di avere introdotto un
supporto per la qualità di servizio che è importante per applicazioni come
quelle multimediali o ``real-time'' che richiedono un qualche grado di
controllo sulla stabilità della banda di trasmissione, sui ritardi o la
dispersione dei temporale del flusso dei pacchetti.


\subsection{Etichette di flusso}
\label{sec:IP_ipv6_flow}
L'introduzione del campo \textit{flow label} può essere usata dall'origine
della comunicazione per etichettare quei pacchetti per i quali si vuole un
trattamento speciale da parte dei router come un una garanzia di banda minima
assicurata o un tempo minimo di instradamento/trasmissione garantito.

Questo aspetto di IPv6 è ancora sperimentale per cui i router che non
supportino queste funzioni devono porre a zero il \textit{flow label} per i
pacchetti da loro originanti e lasciare invariato il campo per quelli in
transito.

Un flusso è una sequenza di pacchetti da una particolare origine a una
particolare destinazione per il quale l'origine desidera un trattamento
speciale da parte dei router che lo manipolano; la natura di questo
trattamento può essere comunicata ai router in vari modi (come un protocollo
di controllo o con opzioni del tipo \textit{hop-by-hop}). 

Ci possono essere più flussi attivi fra un'origine e una destinazione, come
del traffico non assegnato a nessun flusso, un flusso viene identificato
univocamente dagli indirizzi di origine e destinazione e da una etichetta di
flusso diversa da zero, il traffico normale deve avere l'etichetta di flusso
posta a zero.

L'etichetta di flusso è assegnata dal nodo di origine, i valori devono
essere scelti in maniera (pseudo)casuale nel range fra 1 e FFFFFF in modo da
rendere utilizzabile un qualunque sottoinsieme dei bit come chiavi di hash per
i router.

\subsection{Priorità}
\label{sec:IPv6_prio}

Il campo di priorità consente di indicare il livello di priorità dei
pacchetti relativamente agli altri pacchetti provenienti dalla stessa
sorgente. I valori sono divisi in due intervalli, i valori da 0 a 7 sono usati
per specificare la priorità del traffico per il quale la sorgente provvede
un controllo di congestione cioè per il traffico che può essere ``tirato
indietro'' in caso di congestione come quello di TCP, i valori da 8 a 15 sono
usati per i pacchetti che non hanno questa caratteristica, come i pacchetti
``real-time'' inviati a ritmo costante.

Per il traffico con controllo di congestione sono raccomandati i seguenti
valori di priorità a seconda del tipo di applicazione:

\begin{table}[htb]
  \centering
  \footnotesize
  \begin{tabular}{|c|l|}
    \hline
    \textbf{Valore} & \textbf{Tipo di traffico} \\
    \hline
    \hline
    0 & Traffico generico.\\
    1 & Traffico di riempimento (es. news).\\
    2 & Trasferimento dati non interattivo (es. e-mail).\\
    3 & Riservato.\\
    4 & Trasferimento dati interattivo (es. FTP, HTTP, NFS).\\
    5 & Riservato.\\
    \hline
\end{tabular}
\caption{Formato di un indirizzo \textit{site-local}.}
\label{tab:priority}
\end{table}

Per il traffico senza controllo di congestione la priorità più bassa
dovrebbe essere usata per quei pacchetti che si preferisce siano scartati
più facilmente in caso di congestione.


\subsection{Sicurezza a livello IP}
\label{sec:security}

La attuale implementazione di Internet presenta numerosi problemi di
sicurezza, in particolare i dati presenti nelle intestazioni dei vari
protocolli sono assunti essere corretti, il che da adito alla possibilità di
varie tipologie di attacco forgiando pacchetti false, inoltre tutti questi
dati passano in chiaro sulla rete e sono esposti all'osservazione di chiunque
si trovi in mezzo.

Con IPv4 non è possibile realizzare un meccanismo di autenticazione e
riservatezza a un livello inferiore al primo (quello di applicazione), con
IPv6 è stato progettata la possibilità di intervenire al livello di rete (il
terzo) prevedendo due apposite estensioni che possono essere usate per fornire
livelli di sicurezza a seconda degli utenti. La codifica generale di questa
architettura è riportata
nell'\href{http://www.ietf.org/rfc/rfc2401.txt}{RFC~2401}.

Il meccanismo in sostanza si basa su due estensioni:
\begin{itemize}
\item una intestazione di sicurezza (\textit{authentication header}) che
  garantisce al destinatario l'autenticità del pacchetto
\item un carico di sicurezza (\textit{Encrypted Security Payload}) che
  assicura che solo il legittimo ricevente può leggere il pacchetto.
\end{itemize}

Perché tutto questo funzioni le stazioni sorgente e destinazione devono
usare una stessa chiave crittografica e gli stessi algoritmi, l'insieme degli
accordi fra le due stazioni per concordare chiavi e algoritmi usati va sotto
il nome di associazione di sicurezza.

I pacchetti autenticati e crittografati portano un indice dei parametri di
sicurezza (SPI, \textit{Security Parameter Index}) che viene negoziato prima
di ogni comunicazione ed è definito dalla stazione sorgente. Nel caso di
\textit{multicast} dovrà essere lo stesso per tutte le stazioni del gruppo.

\subsection{Autenticazione}
\label{sec:auth} 

Il primo meccanismo di sicurezza è quello dell'intestazione di autenticazione
(\textit{authentication header}) che fornisce l'autenticazione e il controllo
di integrità (ma senza riservatezza) dei pacchetti IP.

L'intestazione di autenticazione ha il formato descritto in
fig.~\ref{fig:autent_estens}: il campo \textit{Next Header} indica
l'intestazione successiva, con gli stessi valori del campo omonimo
nell'intestazione principale di IPv6, il campo \textit{Length} indica la
lunghezza dell'intestazione di autenticazione in numero di parole a 32 bit, il
campo riservato deve essere posto a zero, seguono poi l'indice di sicurezza,
stabilito nella associazione di sicurezza, e un numero di sequenza che la
stazione sorgente deve incrementare di pacchetto in pacchetto.

Completano l'intestazione i dati di autenticazione che contengono un valore di
controllo di integrità (ICV, \textit{Integrity Check Value}), che deve essere
di dimensione pari a un multiplo intero di 32 bit e può contenere un padding
per allineare l'intestazione a 64 bit. Tutti gli algoritmi di autenticazione
devono provvedere questa capacità.

\begin{figure}[!htb]
  \centering
  \includegraphics[width=10cm]{img/ah_option}
    \caption{Formato dell'intestazione dell'estensione di autenticazione.}
    \label{fig:autent_estens}
\end{figure}

L'intestazione di autenticazione può essere impiegata in due modi diverse
modalità: modalità trasporto e modalità tunnel.

La modalità trasporto è utilizzabile solo per comunicazioni fra stazioni
singole che supportino l'autenticazione. In questo caso l'intestazione di
autenticazione è inserita dopo tutte le altre intestazioni di estensione
eccezion fatta per la \textit{Destination Option} che può comparire sia
prima che dopo. 

\begin{figure}[!htb]
  \centering
  \includegraphics[width=10cm]{img/IP_AH_packet}
  \caption{Formato di un pacchetto IPv6 che usa l'opzione di autenticazione.}
  \label{fig:AH_autent_head}
\end{figure}

La modalità tunnel può essere utilizzata sia per comunicazioni fra stazioni
singole che con un gateway di sicurezza; in questa modalità ...


L'intestazione di autenticazione è una intestazione di estensione inserita
dopo l'intestazione principale e prima del carico dei dati. La sua presenza
non ha perciò alcuna influenza sui livelli superiori dei protocolli di
trasmissione come il TCP.


La procedura di autenticazione cerca di garantire l'autenticità del pacchetto
nella massima estensione possibile, ma dato che alcuni campi dell'intestazione
di IP possono variare in maniera impredicibile alla sorgente, il loro valore
non può essere protetto dall'autenticazione.

Il calcolo dei dati di autenticazione viene effettuato alla sorgente su una
versione speciale del pacchetto in cui il numero di salti nell'intestazione
principale è impostato a zero, così come le opzioni che possono essere
modificate nella trasmissione, e l'intestazione di routing (se usata) è posta
ai valori che deve avere all'arrivo.

L'estensione è indipendente dall'algoritmo particolare, e il protocollo è
ancora in fase di definizione; attualmente è stato suggerito l'uso di una
modifica dell'MD5 chiamata \textit{keyed MD5} che combina alla codifica anche
una chiave che viene inserita all'inizio e alla fine degli altri campi.


\subsection{Riservatezza}
\label{sec:ecry}

Per garantire una trasmissione riservata dei dati, è stata previsto la
possibilità di trasmettere pacchetti con i dati criptati: il cosiddetto ESP,
\textit{Encripted Security Payload}. Questo viene realizzato usando con una
apposita opzione che deve essere sempre l'ultima delle intestazioni di
estensione; ad essa segue il carico del pacchetto che viene criptato.

Un pacchetto crittografato pertanto viene ad avere una struttura del tipo di
quella mostrata in fig.~\ref{fig:ESP_criptopack}, tutti i campi sono in chiaro
fino al vettore di inizializzazione, il resto è crittografato.

\begin{figure}[!htb]
  \centering
  \includegraphics[width=10cm]{img/esp_option}
  \caption{Schema di pacchetto crittografato.}
  \label{fig:ESP_criptopack}
\end{figure}



\subsection{Auto-configurazione}
\label{sec:IP_ipv6_autoconf}

Una delle caratteristiche salienti di IPv6 è quella dell'auto-configurazione,
il protocollo infatti fornisce la possibilità ad un nodo di scoprire
automaticamente il suo indirizzo acquisendo i parametri necessari per potersi
connettere a internet.

L'auto-configurazione sfrutta gli indirizzi \textit{link-local}; qualora sul nodo sia
presente una scheda di rete che supporta lo standard IEEE802 (ethernet) questo
garantisce la presenza di un indirizzo fisico a 48 bit unico; pertanto il nodo
può assumere automaticamente senza pericoli di collisione l'indirizzo
\textit{link-local} \texttt{FE80::xxxx:xxxx:xxxx} dove \texttt{xxxx:xxxx:xxxx} è
l'indirizzo hardware della scheda di rete. 

Nel caso in cui non sia presente una scheda che supporta lo standard IEEE802
allora il nodo assumerà ugualmente un indirizzo \textit{link-local} della forma
precedente, ma il valore di \texttt{xxxx:xxxx:xxxx} sarà generato
casualmente; in questo caso la probabilità di collisione è di 1 su 300
milioni. In ogni caso per prevenire questo rischio il nodo invierà un
messaggio ICMP \textit{Solicitation} all'indirizzo scelto attendendo un certo
lasso di tempo; in caso di risposta l'indirizzo è duplicato e il
procedimento dovrà essere ripetuto con un nuovo indirizzo (o interrotto
richiedendo assistenza).

Una volta ottenuto un indirizzo locale valido diventa possibile per il nodo
comunicare con la rete locale; sono pertanto previste due modalità di
auto-configurazione, descritte nelle seguenti sezioni. In ogni caso
l'indirizzo \textit{link-local} resta valido.

\subsection{Auto-configurazione stateless}
\label{sec:stateless}

Questa è la forma più semplice di auto-configurazione, possibile quando
l'indirizzo globale può essere ricavato dall'indirizzo \textit{link-local} cambiando
semplicemente il prefisso a quello assegnato dal provider per ottenere un
indirizzo globale.

La procedura di configurazione è la seguente: all'avvio tutti i nodi IPv6
iniziano si devono aggregare al gruppo di \textit{multicast}
\textit{all-nodes} programmando la propria interfaccia per ricevere i messaggi
dall'indirizzo \textit{multicast} \texttt{FF02::1} (vedi
sez.~\ref{sec:IP_ipv6_multicast}); a questo punto devono inviare un messaggio
ICMP \textit{Router solicitation} a tutti i router locali usando l'indirizzo
\textit{multicast} \texttt{FF02::2} usando come sorgente il proprio indirizzo
\textit{link-local}.

Il router risponderà con un messaggio ICMP \textit{Router Advertisement} che
fornisce il prefisso e la validità nel tempo del medesimo, questo tipo di
messaggio può essere trasmesso anche a intervalli regolari. Il messaggio
contiene anche l'informazione che autorizza un nodo a autocostruire
l'indirizzo, nel qual caso, se il prefisso unito all'indirizzo \textit{link-local} non
supera i 128 bit, la stazione ottiene automaticamente il suo indirizzo
globale.

\subsection{Auto-configurazione stateful}
\label{sec:stateful}

Benché estremamente semplice l'auto-configurazione stateless presenta alcuni
problemi; il primo è che l'uso degli indirizzi delle schede di rete è
molto inefficiente; nel caso in cui ci siano esigenze di creare una gerarchia
strutturata su parecchi livelli possono non restare 48~bit per l'indirizzo
della singola stazione; il secondo problema è di sicurezza, dato che basta
introdurre in una rete una stazione autoconfigurante per ottenere un accesso
legale.

Per questi motivi è previsto anche un protocollo stateful basato su un server
che offra una versione IPv6 del DHCP; un apposito gruppo di \textit{multicast}
\texttt{FF02::1:0} è stato riservato per questi server; in questo caso il nodo
interrogherà il server su questo indirizzo di \textit{multicast} con
l'indirizzo \textit{link-local} e riceverà un indirizzo unicast globale.


\section{Il protocollo ICMP}
\label{sec:ICMP_protocol}

Come già accennato nelle sezioni precedenti, l'\textit{Internet Control
  Message Protocol} è un protocollo di servizio fondamentale per il
funzionamento del livello di rete. Il protocollo ICMP viene trasportato
direttamente su IP, ma proprio per questa sua caratteristica di protocollo di
servizio è da considerarsi a tutti gli effetti appartenente al livello di
rete.

\subsection{L'intestazione di ICMP}
\label{sec:ICMP_header}

Il protocollo ICMP è estremamente semplice, ed il suo unico scopo è quello di
inviare messaggi di controllo; in fig.~\ref{fig:ICMP_header} si è riportata la
struttura dell'intestazione di un pacchetto ICMP generico. 

\begin{figure}[!htb]
  \centering \includegraphics[width=12cm]{img/icmp_head}
  \caption{L'intestazione del protocollo ICMP.}
  \label{fig:ICMP_header}
\end{figure}

Ciascun pacchetto ICMP è contraddistinto dal valore del primo campo, il tipo,
che indica appunto che tipo di messaggio di controllo viene veicolato dal
pacchetto in questione; i valori possibili per questo campo, insieme al
relativo significato, sono riportati in tab.~\ref{tab:ICMP_type}.

\begin{table}[!htb]
  \centering
  \footnotesize
  \begin{tabular}{|l|l|p{9.5cm}|}
    \hline
    \textbf{Valore}&\textbf{Tipo}&\textbf{Significato}\\
    \hline
    \hline
    \texttt{any} & -- & Seleziona tutti i possibili valori \\
    \hline
    \textit{echo-reply}             &0& Inviato in risposta ad un ICMP
                                        \textit{echo-request}.\\ 
    \textit{destination-unreachable}&3& Segnala una destinazione 
                                        irraggiungibile, viene
                                        inviato all'IP sorgente di un
                                        pacchetto quando un router realizza
                                        che questo non può essere inviato a
                                        destinazione.\\
    \textit{source-quench}          &4& Inviato in caso di congestione della
                                        rete per indicare all'IP sorgente di
                                        diminuire il traffico inviato.\\
    \textit{redirect}               &5& Inviato per segnalare un errore di
                                        routing, richiede che la macchina
                                        sorgente rediriga il traffico ad un
                                        altro router da esso specificato.\\
    \textit{echo-request}           &8& Richiede l'invio in risposta di un
                                        \textit{echo-reply}.\\
%    \textit{router-advertisement}   & & \\
%    \textit{router-solicitation}    & & \\
    \textit{time-exceeded}          &11& Inviato quando il TTL di un pacchetto
                                         viene azzerato.\\
    \textit{parameter-problem}      &12& Inviato da un router che rileva dei
                                         problemi con l'intestazione di un
                                         pacchetto.\\
    \textit{timestamp-request}      &13& Richiede l'invio in risposta di un
                                         \textit{timestamp-reply}.\\
    \textit{timestamp-reply}        &14& Inviato in risposta di un
                                         \textit{timestamp-request}.\\
    \textit{info-request}           &15& Richiede l'invio in risposta di un
                                         \textit{info-reply}.\\
    \textit{info-reply}             &16& Inviato in risposta di un
                                         \textit{info-request}.\\
    \textit{address-mask-request}   &17& Richiede l'invio in risposta di un
                                         \textit{address-mask-reply}.\\
    \textit{address-mask-reply}     &18& Inviato in risposta di un
                                         \textit{address-mask-request}.\\
    \hline
  \end{tabular}
  \caption{I valori del \textsl{tipo} per i pacchetti ICMP.}
\label{tab:ICMP_type}
\end{table}

Per alcuni tipi di messaggi ICMP, esiste un secondo campo, detto codice, che
specifica ulteriormente la natura del messaggio; i soli messaggi che
utilizzano un valore per questo campo sono quelli di tipo
\textit{destination-unreachable}, \textit{redirect}, \textit{time-exceeded} e
\textit{parameter-problem}. I possibili valori del codice relativi a ciascuno
di essi sono stati riportati nelle quattro sezioni in cui si è suddivisa
tab.~\ref{tab:ICMP_code}, rispettivamente nell'ordine con cui sono appena
elencati i tipi a cui essi fanno riferimento. 

\begin{table}[!htb]
  \centering
  \footnotesize
  \begin{tabular}{|l|l|}
    \hline
    \textbf{Valore}&\textbf{Codice}\\
    \hline
    \hline
    \textit{network-unreachable}      &0\\
    \textit{host-unreachable}         &1\\
    \textit{protocol-unreachable}     &2\\
    \textit{port-unreachable}         &3 \\
    \textit{fragmentation-needed}     &4\\
    \textit{source-route-failed}      &5\\
    \textit{network-unknown}          &6\\
    \textit{host-unknown}             &7\\
    \textit{host-isolated}            &8\\
    \textit{network-prohibited}       &9\\
    \textit{host-prohibited}          &10 \\
    \textit{TOS-network-unreachable}  &11 \\
    \textit{TOS-host-unreachable}     &12 \\
    \textit{communication-prohibited} &13 \\
    \textit{host-precedence-violation}&14 \\
    \textit{precedence-cutoff}        &15 \\
    \hline
    \textit{network-redirect}         &0  \\
    \textit{host-redirect}            &1  \\
    \textit{TOS-network-redirect}     &2  \\
    \textit{TOS-host-redirect}        &3  \\
    \hline
    \textit{ttl-zero-during-transit}  &0 \\
    \textit{ttl-zero-during-reassembly}&1 \\
    \hline
    \textit{ip-header-bad}            &0 \\
    \textit{required-option-missing}  &1 \\
    \hline
  \end{tabular}
  \caption{Valori del campo \textsl{codice} per il protocollo ICMP.}
\label{tab:ICMP_code}
\end{table}


% LocalWords:  sez Protocol IPv dall' RFC Ethernet Token FDDI Universal host of
% LocalWords:  addressing Best effort l'host router IANA Assigned Number tab to
% LocalWords:  Authority quest'ultime multicast group reserved for CIDR Domain
% LocalWords:  Classless Routing TOS Type Service IPTOS LOWDELAY THROUGHPUT QoS
% LocalWords:  RELIABILITY MINCOST optval anycast unicast fig header version FE
% LocalWords:  priority flow label payload length next hop limit live source FF
% LocalWords:  destination identification fragment checksum TCP UDP ICMPv type
% LocalWords:  service head total fragmentation protocol broadcast broadcasting
% LocalWords:  multicasting path MTU discovery NSAP IPX based geografic local
% LocalWords:  routing format prefix Registry Subscriber Intra Regional SSH
% LocalWords:  Register INTERNIC NCC APNIC subscriber Interface MAC address Reg
% LocalWords:  Subnet Naz Prov Subscr FEBF bootstrap FEC FEFF DNS socket FFFF
% LocalWords:  sull'host loopback scop all nodes routers rip cbt name dhcp HBH
% LocalWords:  agents servers relays solicited extension options route Keyword
% LocalWords:  Authentication Encapsulation ICMP Control Message GGP Gateway ST
% LocalWords:  encapsulation Stream Trasmission Datagram RH FH IDRP ESP Null ip
% LocalWords:  Encrypted Security IGRP OSPF Short First tunnelling FFFFFF hash
% LocalWords:  news FTP NFS authentication Parameter Index ICV Integrity Value
% LocalWords:  padding Option gateway dell'MD keyed Encripted IEEE ethernet any
% LocalWords:  Solicitation netfilter iptables TFTP streaming Normal IGMP
% LocalWords:  stateless solicitation Advertisement stateful Transfer Unit echo
% LocalWords:  l'autoconfigurazione reply request unreachable all'IP quench TTL
% LocalWords:  redirect exceeded parameter problem timestamp info mask port ttl
% LocalWords:  needed failed unknown isolated prohibited communication cutoff
% LocalWords:  precedence violation during reassembly bad required option
% LocalWords:  missing



%%% Local Variables: 
%%% mode: latex
%%% TeX-master: "gapil"
%%% End: 
